\chapter{System Context}
\ChapterLabel{SystemContext}
\label{Section:SystemContextResponsibility}
This chapter defines the mission, system boundary, stakeholders, and external result boundary.
\glsadd{embeddedapplication}\glsadd{controlapp}\glsadd{sample}\glsadd{platformmatrix}\glsadd{conversioncode}\glsadd{countervalue}\glsadd{interfacelogic}\glsadd{qualification}\glsadd{publishedresult}\glsadd{sourcehealth}%


\index{measurement usability}\index{ADC0804}\index{FT245R}\index{Operator}\index{Mission-Integrator}\index{Field-Engineer}
\section{Mission and Boundary}
\label{Section:MeasurementMission}
The mission is to determine whether a measurement can be used. In
Figure~\ref{Figure:SensorContext}, the system of interest includes the
acquisition hardware and the primarily C++ control app on macOS, Windows 11, and iOS. The sensor, ADC0804,
clock/counter, interface logic, and FT245R are external to the control app.
The physical quantity enters the system at the sensor. The Mission application is external to the system. It receives qualified
sensor samples and sensor status through the Open Mission Systems (OMS)
interface. The \ChapterRef{UseCases} boundary is the sensor processing
subsystem, which receives samples from the external Sensor/FTDI hardware.
The arrows show the direction and content of the measurement flow. See
\ArchitectureRef{ARC-103} and \ArchitectureRef{ARC-104} for the allocation
and its requirement links.
\begin{figure}[!htbp]
\centering
\begin{tikzpicture}[
  box/.style={draw=s8sBlue,fill=s8sLightBlue,align=center,
    minimum width=2.9cm,minimum height=0.6cm,font=\small},
  boundary/.style={draw=s8sGray,dashed,inner sep=0.25cm},
  flow/.style={->,>=Latex},
  data/.style={font=\footnotesize,align=left,inner sep=3pt}
]
\node[box,draw=s8sGray,fill=white] (environment) at (-2.6,0) {Measured environment};
\node[box] (sensor) at (-2.6,-1.6) {Sensor};
\node[box] (adc) at (-2.6,-2.8) {ADC0804};
\node[box] (counter) at (2.6,-2.8) {Clock/counter};
\node[box,minimum width=8.1cm] (logic) at (0,-4.0) {Interface logic};
\node[box] (ftdi) at (0,-5.4) {FT245R};
\node[box] (host) at (0,-7.1) {Control app: C++ core\\macOS, Windows 11, iOS\\iOS UI: SwiftUI};
\node[box,draw=s8sGray,fill=white] (consumer) at (0,-10.0) {Mission application};

\coordinate (right edge) at (5.2,-4);
\node[boundary,fit=(sensor)(adc)(counter)(logic)(ftdi)(host)(right edge)] (system) {};
\node[anchor=north east,font=\footnotesize,align=right] at (system.north east)
  {System of interest};

\draw[flow] (environment.south) -- node[data,right]{Physical quantity} (sensor.north);
\draw[flow] (sensor.south) -- node[data,right]{Analog signal} (adc.north);
\draw[flow] (adc.south) -- node[data,right]{8-bit conversion code}
  (adc.south |- logic.north);
\draw[flow] (counter.south) -- node[data,right]{Counter value}
  (counter.south |- logic.north);
\draw[flow] (logic.south) -- node[data,right]{Associated conversion code and\\counter value via parallel FIFO} (ftdi.north);
\draw[flow] (ftdi.south) -- node[data,right]{Associated conversion code and\\counter value via USB} (host.north);
\draw[flow] (host.south) -- node[data,right,pos=0.66]{OMS\\Qualified samples and sensor status} (consumer.north);
\end{tikzpicture}

\caption{Measurement flow across the acquisition hardware, control app, and Mission application boundaries.}
\label{Figure:SensorContext}
\end{figure}

\section{Boundary Information and Open Questions}
\label{Section:ContextBoundaryInformation}
\begin{description}[leftmargin=0pt,labelindent=0pt,labelsep=0pt,style=nextline]
\item[Physical quantity to sensor to ADC0804.]
The sensor responds to the physical quantity and supplies an analog signal
to the ADC0804. Identify the sensor and quantity, then establish the
operating conditions, signal units and range, conversion reference, and
calibration needed to interpret the conversion code
(\hyperref[Question:Q-101]{Q-101}, \hyperref[Question:Q-102]{Q-102}).

\item[ADC0804 and clock/counter to interface logic.]
The ADC0804 supplies an 8-bit conversion code \cite{ADC0804Datasheet}.
The clock/counter supplies a counter value through a separate logical input.
Define the counter width, clock, capture event, reset, rollover, and
relationship to measurement time. Establish how the interface logic associates
the code with the correct counter value and how much capture skew is allowed
(\hyperref[Question:Q-111]{Q-111}, \hyperref[Question:Q-112]{Q-112},
\hyperref[Question:Q-104]{Q-104}).

\item[Interface logic to FT245R.]
The interface logic sends the associated conversion code and counter value
through the FT245R parallel FIFO data path. The FT245R datasheet identifies
the D0--D7 data bus \cite{FT245RDatasheet}. Engineering must establish
electrical compatibility, bus ownership, handshake, transfer order,
record framing, and handling of incomplete transfers and backpressure
(\hyperref[Question:Q-102]{Q-102}, \hyperref[Question:Q-103]{Q-103},
\hyperref[Question:Q-112]{Q-112}).

\item[FT245R to control app.]
USB carries the associated conversion code and counter value to the control
app. Define platform-specific connection access, drivers and APIs for macOS, Windows 11, and iOS, reception and recovery behavior, and
the evidence and decision rules the control app uses to assess measurement
usability. Establish what observations support connection and operating-health
reporting and how the Field-Engineer receives that information
(\hyperref[Question:Q-106]{Q-106}, \hyperref[Question:Q-107]{Q-107},
\hyperref[Question:Q-113]{Q-113}, \hyperref[Question:Q-115]{Q-115}).

\item[Control app to Mission application.]
The control app qualifies measurement usability and supplies qualified samples
and sensor status to the Mission application over OMS. Define the applicable
OMS baseline, field semantics, timing origin and units, and delivery and
error handling with the Mission-Integrator
(\hyperref[Question:Q-114]{Q-114}).
\end{description}

The Operator's decision presentation and any configuration or reverse command
flow remain open in \hyperref[Question:Q-105]{Q-105}.

\section{Stakeholders}
\label{Section:MeasurementStakeholders}
The following roles define the sensor-system stakeholder set. The associated stories distinguish draft needs from the recording and playback needs accepted for inclusion in the story register. Stakeholder roles are distinct from engineering ownership in \hyperref[Table:ProjectOwners]{the owner register}.

\small
\begin{tabularx}{\linewidth}{L{3.6cm}Y}
\caption{Measurement Stakeholder Needs\label{Table:StakeholderNeeds}}\\
\toprule
Stakeholder & Need and story\\\midrule
\endfirsthead
\toprule
Stakeholder & Need and story\\\midrule
\endhead
\bottomrule
\endlastfoot
\phantomsection\label{Stakeholder:Operator}Operator & Needs to see sensor values and their validity, availability, and age so that they can decide whether to use each reading. \UserStoryRef{US-01}.\\
\phantomsection\label{Stakeholder:Acquisition-Developer}Acquisition-Developer & Needs sensor acquisition to continue while prior samples are processed so that mission software receives a continuous stream of readings. \UserStoryRef{US-02}. Accepted diagnostic needs: \UserStoryRef{US-06}, \UserStoryRef{US-08}.\\
\phantomsection\label{Stakeholder:Field-Engineer}Field-Engineer & Needs the sensor subsystem to run on macOS, Windows 11, iOS, and Linux so that they can use any available computer in the field. \UserStoryRef{US-03}. Accepted diagnostic needs: \UserStoryRef{US-05}, \UserStoryRef{US-07}.\\
\phantomsection\label{Stakeholder:Mission-Integrator}Mission-Integrator & Needs a documented interface for sensor data and status so that they can replace the sensor or the mission consumer without redesigning the connection. \UserStoryRef{US-04}. Accepted diagnostic need: \UserStoryRef{US-09}.\\
\end{tabularx}
\normalsize

\section{Stage and Allocation Boundaries}
This stage remains documentation only; no circuit, firmware, host implementation, or product evidence is produced. The control app targets macOS, Windows 11, and iOS under \RequirementRef{REQ-C-103}. \UserStoryRef{US-03} calls for field operation on macOS, Windows 11, iOS, and Linux. The selected platform matrix and connection feasibility remain open in \hyperref[Question:Q-106]{Q-106} and \hyperref[Question:Q-121]{Q-121}; Linux remains an additional stakeholder request. The control-app core remains primarily C++ under \RequirementRef{REQ-C-102}, with an iOS SwiftUI UI under \RequirementRef{REQ-C-104}. The earlier embedded C++ direction remains. The execution target and responsibility of embedded software relative to interface logic remain \hyperref[Question:Q-116]{Q-116}. Selecting ADC0804 and FT245R does not select a processor, FPGA, microcontroller, multiplexer, or host driver. The control app is the host application throughout this book; its name does not imply reverse commands. The iOS UI is selected as SwiftUI; desktop UI technologies and screen-level behavior remain open.
